\documentclass[11pt]{article}
\usepackage[T1]{fontenc}
\usepackage{lmodern}
\usepackage[a4paper,margin=25mm]{geometry}
\usepackage{amsmath,amssymb,amsthm}
\usepackage[hypertexnames=false,hidelinks]{hyperref}
\newcommand{\Z}{\mathbb{Z}}
\newcommand{\V}{\mathsf V}
\newcommand{\E}{\mathsf E}
\newcommand{\loc}[1]{\par\smallskip\noindent{\small\textit{Source: #1}}\par\smallskip}
\setlength{\emergencystretch}{2em}
\begin{document}
\begin{center}
{\large Probability on Trees and Networks\\[2pt]Exercise 6.88 and its Author Comment}\\[4pt]
Russell Lyons\quad Yuval Peres\\[2pt]
Version of 19 August 2026---corrected and updated paperback edition
\end{center}

% BEGIN CONTEXT 6.88
\noindent\textbf{Source and scope (editorial).}
This private reader reproduces selected English author text from the fixed 716-page paperback PDF (21,010,205 bytes). Exercise 6.88 is on physical p.\ 252 (printed p.\ 235), in the collected additional exercises of Chapter 6. Its complete Comment is on physical p.\ 638 (printed p.\ 621). The section the exercise names, Section 6.9, begins on physical p.\ 233 (printed p.\ 216) and defines the anchored expansion constants on physical p.\ 234 (printed p.\ 217); the random subdivision and the exponential-tail condition are defined in the body on physical p.\ 238 (printed p.\ 221), which also prints the constant $\Phi^{**}_{\E}$ and Theorem 6.50. The edge-set notation that the definition of $\Phi^{**}_{\E}$ invokes, equation (6.33), is printed on physical p.\ 217 (printed p.\ 200). The one place in the body that cites the exercise is the necessity remark following the proof of Theorem 6.50 on physical p.\ 240 (printed p.\ 223). All of this is reprinted below as uncounted context.

The separate errata sheet dated 19 August 2026 contains exactly one Chapter 6 entry, the p.\ 616 correction to Exercise 6.34. Its only other entry touching Chapter 6 material is on p.\ 650, where two bibliographic items dated 2015 that are printed on pp.\ 509 and 510 are replaced. No entry touches physical pp.\ 217, 233, 234, 238, 240, 252 or 638. Accordingly no revision is applied anywhere in this reader, and no printed text here is altered. No source mathematics is repaired or independently certified.

\medskip\noindent\textbf{Literal source features kept as printed (no synchronization, no repair).}
\begin{itemize}
\item The anchored-expansion constant is printed with a single star on physical p.\ 252, in the question, and with a single star in the closing display of the Comment on physical p.\ 638. A direct measurement of the glyph geometry of both pages (recorded in this batch's instrument receipts) finds one raised script-size asterisk per occurrence, at 6.30\,pt, positioned after the capital $\Phi$; the PDF text layer of physical p.\ 252 yields three occurrences of $\Phi$ followed by one asterisk and none followed by two, and physical p.\ 638 yields one occurrence with a single asterisk and none with two. So the asked constant and the constant that the Comment closes with are the same printed symbol.
\item The constant with a double star is a different printed symbol and appears only on physical pp.\ 238 and 240, where it is set as two separate raised script-size asterisks at 7.35\,pt; the text layer of both pages yields six occurrences of the double form and no single-star substitution. It is transcribed here as $\Phi^{**}_{\E}$ exactly where the body prints it, and the single-star form $\Phi^{*}_{\E}$ exactly where the body prints that.
\item The unanchored expansion constant of Exercise 6.23 on physical p.\ 238 carries no star at all in print; no asterisk glyph sits after its $\Phi$. It is transcribed as $\Phi_{\E}$.
\item The exponents of the Comment on physical p.\ 638 are true raised superscripts, not baseline subscripts: the $n$ of $2^{n}$, the $\epsilon$ and $n$ of $e^{-\epsilon n}$, the $2n$ in the exponent of $(1-e^{-\epsilon n})^{2n}$, and the further raised $2^{n}$ there, which is measured at 5.25\,pt, one level above the 6.30\,pt script size of the exponents that carry it. The last is transcribed as printed, preserving the two levels. The subscripts $i=1$ under the two sums, the subscript $1$ of $L_{1}$, and the subscript $\E$ of $\partial_{\E}$ and of $\Phi^{*}_{\E}$ are set on the baseline and are ordinary subscripts.
\item The Comment prints $L_{e}$ and $\sum_{e\in E(S)}L_{e}$ over the edge set $E(S)$, and $\partial_{\E}S$ over the vertex set $S$; the denominator is neither $\sum_{e\in E^{*}(S)}L_{e}$ nor $|E(S)|$, and it is not synchronized with the corresponding display of the Theorem 6.50 proof on physical p.\ 240, which uses $E^{*}(S)$. Both are transcribed as printed.
\item The Comment and the body set the law and its members in angle brackets, $\langle L_{i}\rangle$ and $\langle L_{e}\rangle_{e\in\E(G)}$; these are transcribed as printed.
\item The space character after the symbol $\nu$ is missing from the PDF text layer of every page of this exercise, not only from the Comment: the extracted text reads ``$\nu$has'' on physical p.\ 252, ``$\nu$has'' on physical p.\ 238 and ``$\nu$of'' on physical p.\ 638, and the other occurrences of $\nu$ on those pages behave the same way. The printed spacing is nevertheless ordinary there: measured gap from the right edge of the $\nu$ span to the left edge of the following word is 3.04\,pt at 9.11\,pt type on p.\ 252, 3.23\,pt and 3.64\,pt at 10.13\,pt on p.\ 238, and 2.63\,pt at 8.10\,pt on p.\ 638. No difference of spacing between the Comment and the body is claimed, and no space is added to or removed from the transcription on that ground.
\item The stretched delimiter glyphs of some formulas do not extract as delimiter characters: the text layer of the fixed version writes them as control bytes of the font \texttt{txex-bar}, and the byte value does not identify which delimiter is printed. Every line of all seven source pages was scanned. Physical p.\ 217 carries two such delimiter pairs in the two displayed definitions of $E(K)$ and $E^{*}(K)$ and one brace pair in $G\!\restriction\!K=\{K,E(K)\}$; p.\ 233 carries three parenthesis pairs around $|W_{k}|^{\pi}$ in $2/\phi(\,\cdot\,)$; p.\ 234 carries three parenthesis pairs of $1+\Phi^{*}_{\E}$; p.\ 240 carries two parenthesis pairs of $c(1+\Phi^{**}_{\E})$ and two of the probabilities $P(A(n,h))$ and $P(n\le|\V(T)|<\infty)$; p.\ 252 carries two parenthesis pairs of $1+|K|$, two bar glyphs of the norm $|B|$ in Exercise 6.80(a), and four bar glyphs of the two norms in Exercise 6.86. Physical pp.\ 238 and 638 carry no control byte at all: the angle brackets $\langle L_{e}\rangle_{e\in\E(G)}$ and $\langle L_{i}\rangle$ of the law, the brace pair of $|\partial_{\E}K|$ that follows, and the brackets of the Comment's expectations extract as ordinary characters. All of this is an extraction artefact of the delimiters, not a printed character, and nothing is transcribed from it. The delimiters of this reader are ordinary characters.
\item In the displayed bound of the Comment the numerator is printed as $2n+1$ and the denominator as $cn$, and the following relation is the two-line approximation symbol, $\approx$, not an equality. Both are transcribed as printed.
\item The reflowed line breaks of the question and of the Comment are not content; the sentences are printed as continuous prose in the source and are set here as continuous prose.
\end{itemize}

\medskip\noindent\textbf{Use of the exercise in the body of the work.}
The number 6.88 occurs on exactly three physical pages of the fixed version: 240, 252 and 638. The occurrence on physical p.\ 240 is the author's necessity remark, ``The exponential tail condition is necessary to ensure the positivity of $\Phi^{*}_{\E}(G_{\nu})$; see Exercise 6.88.'', which sits directly after the proof of Theorem 6.50 and is the only in-body citation of the exercise. There is no theorem, proposition, lemma or example in the fixed version that states this exercise's conclusion under its own name. The heading ``Chapter 7'' that the text layer of physical p.\ 638 reports after the Comment begins at a vertical position below the Comment's last line and belongs to the collected Comments, not to the Comment on 6.88; nothing after the closing probabilistic statement of the Comment is reprinted here as part of it.

\subsection*{Section 6.9, Anchored Expansion and Percolation}
\loc{Physical p.\ 233 (printed p.\ 216).}
\textbf{6.9 Anchored Expansion and Percolation.} Recall the probability measure $P_{p}$ defining Bernoulli percolation from Section 5.2 and the critical probability $p_{c}$. There, we took a random subset of edges, but an alternative is to take the subgraph induced by a random subset of vertices. This alternative is called site percolation. The adjective Bernoulli applies when each vertex is present independently with the same probability. When we need different notation for these two processes, we use $P^{\mathrm{site}}_{p}$ and $P^{\mathrm{bond}}_{p}$ for the two product measures on $2^{\V}$ and $2^{\E}$, respectively, and $p^{\mathrm{site}}_{c}$ and $p^{\mathrm{bond}}_{c}$ for the two critical probabilities. If we don't indicate whether the percolation is bond or site and both make sense in context, then results we state should be taken to apply to both types of percolation.

Grimmett, Kesten, and Zhang (1993) showed that simple random walk on the infinite cluster of Bernoulli percolation in $\Z^{d}$ (when $p>p_{c}$) is transient for $d\ge3$; in other words, in Euclidean lattices, transience is preserved when the whole lattice is replaced by an infinite percolation cluster. (We gave a proof of part of this result in Section 5.5, where we also noted that another proof based on ideas in the present section was given by Pete (2008).) In general, what aspects of Cayley graphs are preserved under percolation?

\textbf{Conjecture 6.44.} (Percolation and Transience) \textit{If $G$ is a transient Cayley graph, then a.s.\ every infinite cluster of Bernoulli percolation on $G$ is transient.}

\textbf{Conjecture 6.45.} (Percolation and Speed) \textit{Let $G$ be a Cayley graph. Then simple random walk on $G$ has positive speed iff simple random walk on infinite clusters of Bernoulli percolation has positive speed a.s.}

These conjectures were made by Benjamini, Lyons, and Schramm (1999), who proved that simple random walk on an infinite cluster of any nonamenable Cayley graph has positive speed. One might hope to use Proposition 6.9 to establish this result, but, in fact, the infinite clusters are amenable:

\subsection*{The anchored expansion constants}
\loc{Physical p.\ 234 (printed p.\ 217).}
\textbf{Exercise 6.19.} For any $p<1$, every infinite cluster $K$ of Bernoulli($p$) percolation on any graph $G$ of bounded degree has $\Phi_{\E}(K)=0$ a.s.

On the other hand, the infinite clusters might satisfy the following weaker ``anchored expansion'' property, which is known to imply positive speed (Theorem 6.53).

Fix $o\in\V(G)$. The \textit{anchored expansion constants} of $G$ are
\[
\Phi^{*}_{\E}(G):=\lim_{n\to\infty}\inf\Bigl\{\frac{|\partial_{\E}K|}{|K|}\ ;\ o\in K\subset\V,\ G\!\restriction\!K\text{ is connected},\ n\le|K|<\infty\Bigr\}
\]
and
\[
\Phi^{*}_{\V}(G):=\lim_{n\to\infty}\inf\Bigl\{\frac{|\partial_{\V}K|}{|K|}\ ;\ o\in K\subset\V,\ G\!\restriction\!K\text{ is connected},\ n\le|K|<\infty\Bigr\}.
\]
These are closely related to the number $\psi(G,o,0)$ of (6.57), but have the advantage that $\Phi^{*}_{\E}(G)$ and $\Phi^{*}_{\V}(G)$ do not depend on the choice of the basepoint $o$. We say that a graph $G$ has \textit{anchored expansion} if $\Phi^{*}_{\E}(G)>0$.

\subsection*{Random subdivision and the exponential tail}
\loc{Physical p.\ 238 (printed p.\ 221).}
Let $G$ be an infinite graph of bounded degree, and pick a probability distribution $\nu$ on the positive integers. Replace each edge $e\in\E(G)$ by a path of $L_{e}\ge1$ edges, where the random variables $\langle L_{e}\rangle_{e\in\E(G)}$ are independent with law $\nu$. Let $G_{\nu}$ denote the random graph obtained from $G$ in this way. We call $G_{\nu}$ a \textit{random subdivision} of $G$. Say that $\nu$ has an \textit{exponential tail} if, for some $\epsilon>0$ and all sufficiently large $\ell$, we have $\nu[\ell,\infty)<e^{-\epsilon\ell}$. This is equivalent to the condition that if $X\sim\nu$, then $\E[s^{X}]<\infty$ for some $s>1$.

\textbf{Exercise 6.23.} Show that if the support of $\nu$ is unbounded, then $\Phi_{\E}(G_{\nu})=0$ a.s.

\subsection*{The notation of (6.33) used by the definition of $\Phi^{**}_{\E}$}
\loc{Physical p.\ 217 (printed p.\ 200).}
For $K\subseteq\V$, define
\[
E(K):=\overline{[x,y]\in\E\,;\,x,y\in K}\qquad(6.33)
\]
and
\[
E^{*}(K):=\overline{[x,y]\in\E\,;\,x\in K\text{ or }y\in K}.
\]
Thus, $\partial_{\E}K=E^{*}(K)\setminus E(K)$ and the graph induced by $K$ is $G\!\restriction\!K=\overline{K,E(K)}$.

\subsection*{The constant $\Phi^{**}_{\E}$ and Theorem 6.50}
\loc{Physical p.\ 238 (printed p.\ 221).}
Define another anchored expansion constant,
\[
\Phi^{**}_{\E}(G):=\lim_{n\to\infty}\inf\Bigl\{\frac{|\partial_{\E}K|}{|E(K)|}\ ;\ o\in K\subset\V,\ G\!\restriction\!K\text{ is connected},\ n\le|K|<\infty\Bigr\},
\]
with notation as in (6.33). Since $|E(K)|\ge|K|-1$, we have
\[
\Phi^{**}_{\E}(G)\le\Phi^{*}_{\E}(G)
\]
for every graph $G$. Conversely, if the maximum degree of $G$ is $D$, then
\[
\Phi^{**}_{\E}(G)\ge\Phi^{*}_{\E}(G)/D
\]
since $|E(K)|\le D|K|$. On the other hand, for trees $G$, we have $|E(K)|=|K|-1$, so
\[
\Phi^{**}_{\E}(G)=\Phi^{*}_{\E}(G)\text{ for trees.}
\]

\textbf{Theorem 6.50.} (Anchored Expansion and Subdivision) \textit{Suppose that $\Phi^{**}_{\E}(G)>0$. If $\nu$ has an exponential tail, then the random subdivision satisfies $\Phi^{**}_{\E}(G_{\nu})>0$ a.s. In particular, if $G$ has bounded degree and $\Phi^{*}_{\E}(G)>0$, then $\Phi^{*}_{\E}(G_{\nu})>0$ a.s.}

\subsection*{The necessity remark and the Galton--Watson corroboration}
\loc{Physical p.\ 240 (printed p.\ 223).}
The proof of Theorem 6.50 closes on this page with
\[
\text{so we conclude that}\quad \Phi^{**}_{\E}(G_{\nu})\ \ge\ \Phi^{**}_{\E}(G)\big/\Bigl(c\bigl(1+\Phi^{**}_{\E}(G)\bigr)\Bigr)\ >\ 0.\quad\blacktriangleleft
\]
It is followed immediately by: The exponential tail condition is necessary to ensure the positivity of $\Phi^{*}_{\E}(G_{\nu})$; see Exercise 6.88.

The same page then argues that Galton--Watson trees have anchored expansion, and in Case (ii) of the proof of Theorem 6.52 uses this exercise's mechanism from the other side. For a supercritical Galton--Watson tree $T$, given nonextinction, $\Phi^{*}_{\E}(T)>0$ a.s. In Case (ii), $p_{0}=0$ and $p_{1}>0$: let $x$ be the vertex closest to the root that has at least two children; then $T^{x}$ has the law of a random subdivision $G_{\nu}$ of another Galton--Watson tree $G$, and $T$ differs from $T^{x}$ by a finite path. Here $G$ is generated according to the offspring distribution $\langle p'_{k}\,;\,k\ge0\rangle$, where $p'_{k}:=p_{k}/(1-p_{1})$ for $k=2,3,\dots$ and $p'_{0}:=p'_{1}:=0$, and $\nu$ is the geometric distribution with parameter $1-p_{1}$. By Theorem 6.50 and the fact that $\Phi^{**}_{\E}=\Phi^{*}_{\E}$ for trees, $\Phi^{*}_{\E}(T)=\Phi^{*}_{\E}(T^{x})=\Phi^{*}_{\E}(G_{\nu})>0$ a.s.
% END CONTEXT 6.88

\clearpage
% BEGIN PRIMARY STATEMENT 6.88
\section*{Exercise 6.88}
\loc{Physical p.\ 252 (printed p.\ 235).}
\begin{sloppypar}
\noindent\textbf{Exercise 6.88.} Let $G$ be a binary tree. Show that if $\nu$ has a tail that decays slower than exponentially, then $\Phi^{*}_{\E}(G)>0$ yet $\Phi^{*}_{\E}(G_{\nu})=0$ a.s., where $G_{\nu}$ is the random subdivision defined in Section 6.9.
\end{sloppypar}
% END PRIMARY STATEMENT 6.88

% BEGIN PRIMARY PROOF 6.88
\section*{Complete author Comment on Exercise 6.88}
\loc{Physical p.\ 638 (printed p.\ 621).}
\textbf{6.88.} Suppose that the distribution $\nu$ of $L$ does not have an exponential tail. Then, for every $c>0$ and every $\epsilon>0$, we have
\[
\mathsf P\Bigl[\sum_{i=1}^{n}L_{i}\ge cn\Bigr]\ \ge\ \mathsf P[L_{1}\ge cn]\ \ge\ e^{-\epsilon n}
\]
for infinitely many $n$s, where $\langle L_{i}\rangle$ are i.i.d.\ with law $\nu$. Let $G$ be a binary tree with the root $o$ as the basepoint. Pick a collection of $2^{n}$ pairwise disjoint paths from level $n$ to level $2n$. Then
\begin{multline*}
\mathsf P\Bigl[\text{along at least one of these }2^{n}\text{ paths }\sum_{i=1}^{n}L_{i}\ge cn\Bigr]\\
\ge\ 1-(1-e^{-\epsilon n})^{2^{n}}\ \ge\ 1-\exp\bigl(-e^{-\epsilon n}2^{n}\bigr)\ \to\ 1.
\end{multline*}
With probability very close to 1 (depending on $n$), there is a path from level $n$ to $2n$ along which $\sum_{i=1}^{n}L_{i}\ge cn$. Take such a path and extend it to the root, $o$. Let $S$ be the set of vertices in the extended path from the root $o$ to level $2n$. Then
\[
\frac{|\partial_{\E}S|}{\sum_{e\in E(S)}L_{e}}\ \le\ \frac{2n+1}{cn}\ \approx\ \frac{2}{c}.
\]
Since $c$ can be arbitrarily large, $\Phi^{*}_{\E}(G_{\nu})=0$ a.s.
% END PRIMARY PROOF 6.88

\medskip\noindent\textbf{Editorial note on what this reader does not claim.}
The Comment above is reproduced in full and is the author's argument for the exercise as stated. This reader adds no step to it, repairs no step of it, and certifies no step of it. In particular, the passage from ``$\nu$ does not have an exponential tail'' to the two inequalities displayed in the Comment's second sentence is printed as an assertion and is not derived on either the statement page or the Comment page; it is left exactly as printed. The relation between the two normalizations that appear in the printed record is likewise left as printed: the Comment's final display bounds $|\partial_{\E}S|$ over $\sum_{e\in E(S)}L_{e}$, while the definition of $\Phi^{*}_{\E}$ on physical p.\ 234 bounds $|\partial_{\E}K|$ over $|K|$, and the two are not identified here. The statement of the exercise and the closing display of the Comment both carry the single-star constant, as measured and recorded above.
\end{document}
